\medskip

Un biologiste a modélisé l'évolution d'une population de bactéries (en milliers d'entités) par la fonction $f$ définie sur $[0~;~+\infty[$ par

\[f(t) = \text{e}^3 - \text{e}^{-0,5t^2 + t + 2}\]

où $t$ désigne le temps en heures depuis le début de l'expérience.

\medskip

À partir de cette modélisation, il propose les trois affirmations ci-dessous.

Pour chacune d'elles, indiquer, en justifiant, si elle est vraie ou fausse.

\begin{list}{\textbullet}{}
\item\textbf{Affirmation 1} : \og La population augmente en permanence \fg.

\item\textbf{Affirmation 2} : \og  À très long terme, la population dépassera \np{21000} bactéries\fg.

\item\textbf{Affirmation 3} : \og La population de bactéries aura un effectif de \np{10000} à deux reprises au cours du temps \fg.
\end{list}
